\documentclass[10pt]{article}
\usepackage[T1]{fontenc}
\usepackage[utf8]{inputenc}
\usepackage[french]{babel}
\usepackage{amsmath,amssymb,mathrsfs,mathpazo}
\usepackage{geometry}
\geometry{a4paper,margin=2cm}
\usepackage{graphicx,array,tabularx,tabularray,colortbl,multirow,multicol,booktabs,enumitem,pifont,ulem,fancybox,diagbox,cancel}
\usepackage{pstricks,pstricks-add,pst-plot,pst-tree,pst-node,pst-text,pst-func,pst-math,pst-bspline,pst-3dplot,pst-eucl,pst-solides3d}
\let\pscalculate\undefined
\usepackage{pst-bezier}
\usepackage{tikz,pgfplots,tkz-tab}
\usetikzlibrary{arrows,arrows.meta,patterns,calc,positioning,trees,decorations.markings,shapes.geometric,shapes.misc,shapes.arrows,intersections,decorations.pathmorphing,decorations.pathreplacing,backgrounds,shadows,babel}
\pgfplotsset{compat=1.18}
\usepgfplotslibrary{fillbetween}
\usepackage[np]{numprint}
\usepackage{listings,makecell,pas-tableur,twoopt}
\let\barre\undefined
\usepackage[pstricks]{tablvar}
\definecolor{darkgreen}{RGB}{0,100,0}
\providecommand{\red}{\color{red}}
\providecommand{\blue}{\color{blue}}
\providecommand{\green}{\color{green}}
\providecommand{\black}{\color{black}}
\providecommand{\white}{\color{white}}
\providecommand{\gray}{\color{gray}}
\providecommand{\cyan}{\color{cyan}}
\providecommand{\magenta}{\color{magenta}}
\providecommand{\yellow}{\color{yellow}}
\providecommand{\violet}{\color{violet}}
\providecommand{\orange}{\color{orange}}
\providecommand{\euro}{\texteuro}
\providecommand{\EUR}{\texteuro}
\providecommand{\ieme}{\textsuperscript{e}}
\providecommand{\degres}{\ensuremath{^\circ}}
\providecommand{\degree}{\ensuremath{^\circ}}
\providecommand{\vv}[1]{\overrightarrow{#1}}
\providecommand{\Coord}[1]{#1}
\providecommand{\decosix}{}
\providecommand{\starredbullet}{\ensuremath{\bullet}}
\pagestyle{empty}
\setlength{\parindent}{0pt}
\usepackage{adjustbox}\DeclareUnicodeCharacter{2236}{:}\DeclareUnicodeCharacter{2190}{\ensuremath{\leftarrow}}
\providecommand{\R}{\mathbb{R}}
\providecommand{\N}{\mathbb{N}}
\providecommand{\Z}{\mathbb{Z}}
\providecommand{\Q}{\mathbb{Q}}
\providecommand{\D}{\mathbb{D}}
\providecommand{\C}{\mathbb{C}}
\providecommand{\e}{{\mathrm{e}}}
\providecommand{\vect}[1]{\overrightarrow{#1}}
\providecommand{\vectt}[1]{\overrightarrow{\mathrm{#1}}}
\providecommand{\vvt}[1]{\overrightarrow{\mathrm{#1}}}
\providecommand{\np}[2][]{{#2~#1}}
\providecommand{\vv}[1]{{\overrightarrow{#1}}}
\renewcommand{\d}{{\mathrm{d}}}
\providecommand{\degres}{{^\circ}}
\providecommand{\degre}{{^\circ}}
\providecommand{\blue}{\color{blue}}
\providecommand{\red}{\color{red}}
\providecommand{\enskip}{\quad}
\providecommand{\llbracket}{\lbrack\!\lbrack}
\providecommand{\rrbracket}{\rbrack\!\rbrack}
\providecommand{\euro}{{\text{€}}}
\providecommand{\no}{{\text{n°}}}
\providecommand{\interval}[2]{{[#1;#2]}}
\providecommand{\intervFO}[2]{{[#1;#2[}}
\providecommand{\dotfill}{\ldots}

\providecommand{\ds}{\displaystyle}
\providecommand{\textding}[1]{#1}
\providecommand{\Oij}{\(\mathrm{O};\vec{\imath},\vec{\jmath}\)}
\providecommand{\Oijk}{\(\mathrm{O};\vec{\imath},\vec{\jmath},\vec{k}\)}
\providecommand{\Ouv}{\(\mathrm{O};\vec{u},\vec{v}\)}
\providecommand{\intervFF}[2]{\ensuremath{[#1;#2]}}
\renewcommand{\euro}{\ifmmode\text{\texteuro}\else\texteuro\fi}
\DeclareUnicodeCharacter{00B0}{\ifmmode^{\circ}\else\textdegree\fi}
\DeclareUnicodeCharacter{20AC}{\ifmmode\text{\texteuro}\else\texteuro\fi}
% Énoncés : Éducation nationale. Transcription et corrigés : APMEP, Denis Vergès et contributeurs. Découpage et classement : CoopMaths, Sébastien Lozano, Éric Elter et contributeurs. Adaptation de la présentation : Maths Explorer.
% CC BY-SA : https://coopmaths.fr/www/about/
% https://coopmaths.fr/alea/static/bac/2021/tex/bac_2021_03_sujet1_1.tex
\begin{document}
\noindent\textit{APMEP / CoopMaths. CC BY-SA. Présentation adaptée par Maths Explorer.}\par\medskip
\medskip

Dans une école de statistique, après étude des dossiers des candidats, le recrutement se fait de deux façons :

\setlength\parindent{1cm}
\begin{itemize}
\item[$\bullet~~$] 10\,\% des candidats sont sélectionnés sur dossier. Ces candidats doivent ensuite passer un oral à l'issue duquel 60\,\% d'entre eux sont finalement admis à l'école.
\item[$\bullet~~$] Les candidats n'ayant pas été sélectionnés sur dossier passent une épreuve écrite à l'issue de laquelle 20\,\% d'entre eux sont admis à l'école.
\end{itemize}
\setlength\parindent{0cm}

\bigskip

\textbf{Partie 1}

\medskip

On choisit au hasard un candidat à ce concours de recrutement. On notera:

\setlength\parindent{1cm}
\begin{itemize}
\item[$\bullet~~$] $D$ l'évènement \og le candidat a été sélectionné sur dossier \fg{} ;
\item[$\bullet~~$] $A$ l'évènement \og le candidat a été admis à l'école \fg{} ;
\item[$\bullet~~$] $\overline{D}$ et $\overline{A}$ les évènements contraires des évènements $D$ et $A$ respectivement.
\end{itemize}
\setlength\parindent{0cm}

\medskip

\begin{enumerate}
\item Traduire la situation par un arbre pondéré.
\item Calculer la probabilité que le candidat soit sélectionné sur dossier et admis à l'école.
\item Montrer que la probabilité de l'évènement  $A$ est égale à $0,24$.
\item On choisit au hasard un candidat admis à l'école. Quelle est la probabilité que son dossier n'ait pas été sélectionné?
\end{enumerate}

\bigskip

\textbf{Partie 2}

\medskip

\begin{enumerate}
\item On admet que la probabilité pour un candidat d'être admis à l'école est égale à $0,24$.

On considère un échantillon de sept candidats choisis au hasard, en assimilant ce choix à un tirage au sort avec remise. On désigne par $X$ la variable aléatoire dénombrant les candidats admis à l'école parmi les sept tirés au sort.
	\begin{enumerate}
		\item On admet que la variable aléatoire $X$ suit une loi binomiale. Quels sont les paramètres de cette loi?
		\item Calculer la probabilité qu'un seul des sept candidats tirés au sort soit admis à l'école. On donnera une réponse arrondie au centième.
		\item Calculer la probabilité qu'au moins deux des sept candidats tirés au sort soient admis à cette école. On donnera une réponse arrondie au centième.
	\end{enumerate}
\item  Un lycée présente $n$ candidats au recrutement dans cette école, où $n$ est un entier naturel non nul.

On admet que la probabilité pour un candidat quelconque du lycée d'être admis à l'école est égale à $0,24$ et que les résultats des candidats sont indépendants les uns des autres.
	\begin{enumerate}
		\item Donner l'expression, en fonction de $n$, de la probabilité qu'aucun candidat issu de ce lycée ne soit admis à l'école.
		\item À partir de quelle valeur de l'entier $n$ la probabilité qu'au moins un élève de ce lycée soit admis à l'école est-elle supérieure ou égale à $0,99$ ?
	\end{enumerate}
\end{enumerate}

\bigskip

\end{document}
